% Диаграмма сравнения производительности: эмулятор (без/со профайлингом) vs V100.
% Данные из таблицы tab:integration_results (ch8_testing.tex).
% Лог. шкала по Y: значения V100 (мкс) и эмулятора (мс) различаются на ~5 порядков.
% Подключается в главе через: % Диаграмма сравнения производительности: эмулятор (без/со профайлингом) vs V100.
% Данные из таблицы tab:integration_results (ch8_testing.tex).
% Лог. шкала по Y: значения V100 (мкс) и эмулятора (мс) различаются на ~5 порядков.
% Подключается в главе через: % Диаграмма сравнения производительности: эмулятор (без/со профайлингом) vs V100.
% Данные из таблицы tab:integration_results (ch8_testing.tex).
% Лог. шкала по Y: значения V100 (мкс) и эмулятора (мс) различаются на ~5 порядков.
% Подключается в главе через: % Диаграмма сравнения производительности: эмулятор (без/со профайлингом) vs V100.
% Данные из таблицы tab:integration_results (ch8_testing.tex).
% Лог. шкала по Y: значения V100 (мкс) и эмулятора (мс) различаются на ~5 порядков.
% Подключается в главе через: \input{images/perf_comparison.tex}
% Требует в преамбуле: \usepackage{pgfplots}\pgfplotsset{compat=1.18}
\begin{figure}[h]
  \centering
  \begin{tikzpicture}
    \begin{axis}[
        width=15cm, height=8.5cm,
        ybar=1.5pt,
        bar width=8pt,
        ymode=log,
        log basis y=10,
        log origin y=infty,
        ymin=0.001, ymax=10000,
        ylabel={Время исполнения, мс (лог. шкала)},
        symbolic x coords={vadd,gelu,softmax,fft,mmul,attention,conj\_grad},
        xtick=data,
        x tick label style={font=\small},
        ymajorgrids=true,
        grid style={gray!25},
        enlarge x limits=0.08,
        legend style={at={(0.5,-0.16)}, anchor=north, legend columns=3,
                      font=\small, draw=none},
        legend cell align=left,
      ]
      % Эмулятор без профайлинга
      \addplot[ybar, draw=black, line width=0.4pt, fill=gray!50] coordinates {
        (vadd,20.5) (gelu,46.8) (softmax,60.9) (fft,52.6)
        (mmul,111.7) (attention,300.5) (conj\_grad,890.6)};
      % Эмулятор со сбором профайлинга
      \addplot[ybar, draw=black, line width=0.4pt, fill=gray!90] coordinates {
        (vadd,34.8) (gelu,121.8) (softmax,268.4) (fft,218.0)
        (mmul,970.5) (attention,2074.1) (conj\_grad,1879.7)};
      % V100-PCIE (значения переведены из мкс в мс)
      \addplot[ybar, draw=black, line width=0.4pt, fill=white] coordinates {
        (vadd,0.0072) (gelu,0.0080) (softmax,0.0078) (fft,0.0089)
        (mmul,0.0178) (attention,0.1319) (conj\_grad,0.4763)};
      \legend{Эмулятор (без проф.), Эмулятор (со сбором проф.), V100-PCIE}
    \end{axis}
  \end{tikzpicture}
  \caption{Время исполнения ядер: эмулятор без профайлинга, со сбором
           профайлинга и реальный GPU V100-PCIE. Логарифмическая шкала;
           значения V100 переведены из микросекунд в миллисекунды.}
  \label{fig:perf_comparison}
\end{figure}

% Требует в преамбуле: \usepackage{pgfplots}\pgfplotsset{compat=1.18}
\begin{figure}[h]
  \centering
  \begin{tikzpicture}
    \begin{axis}[
        width=15cm, height=8.5cm,
        ybar=1.5pt,
        bar width=8pt,
        ymode=log,
        log basis y=10,
        log origin y=infty,
        ymin=0.001, ymax=10000,
        ylabel={Время исполнения, мс (лог. шкала)},
        symbolic x coords={vadd,gelu,softmax,fft,mmul,attention,conj\_grad},
        xtick=data,
        x tick label style={font=\small},
        ymajorgrids=true,
        grid style={gray!25},
        enlarge x limits=0.08,
        legend style={at={(0.5,-0.16)}, anchor=north, legend columns=3,
                      font=\small, draw=none},
        legend cell align=left,
      ]
      % Эмулятор без профайлинга
      \addplot[ybar, draw=black, line width=0.4pt, fill=gray!50] coordinates {
        (vadd,20.5) (gelu,46.8) (softmax,60.9) (fft,52.6)
        (mmul,111.7) (attention,300.5) (conj\_grad,890.6)};
      % Эмулятор со сбором профайлинга
      \addplot[ybar, draw=black, line width=0.4pt, fill=gray!90] coordinates {
        (vadd,34.8) (gelu,121.8) (softmax,268.4) (fft,218.0)
        (mmul,970.5) (attention,2074.1) (conj\_grad,1879.7)};
      % V100-PCIE (значения переведены из мкс в мс)
      \addplot[ybar, draw=black, line width=0.4pt, fill=white] coordinates {
        (vadd,0.0072) (gelu,0.0080) (softmax,0.0078) (fft,0.0089)
        (mmul,0.0178) (attention,0.1319) (conj\_grad,0.4763)};
      \legend{Эмулятор (без проф.), Эмулятор (со сбором проф.), V100-PCIE}
    \end{axis}
  \end{tikzpicture}
  \caption{Время исполнения ядер: эмулятор без профайлинга, со сбором
           профайлинга и реальный GPU V100-PCIE. Логарифмическая шкала;
           значения V100 переведены из микросекунд в миллисекунды.}
  \label{fig:perf_comparison}
\end{figure}

% Требует в преамбуле: \usepackage{pgfplots}\pgfplotsset{compat=1.18}
\begin{figure}[h]
  \centering
  \begin{tikzpicture}
    \begin{axis}[
        width=15cm, height=8.5cm,
        ybar=1.5pt,
        bar width=8pt,
        ymode=log,
        log basis y=10,
        log origin y=infty,
        ymin=0.001, ymax=10000,
        ylabel={Время исполнения, мс (лог. шкала)},
        symbolic x coords={vadd,gelu,softmax,fft,mmul,attention,conj\_grad},
        xtick=data,
        x tick label style={font=\small},
        ymajorgrids=true,
        grid style={gray!25},
        enlarge x limits=0.08,
        legend style={at={(0.5,-0.16)}, anchor=north, legend columns=3,
                      font=\small, draw=none},
        legend cell align=left,
      ]
      % Эмулятор без профайлинга
      \addplot[ybar, draw=black, line width=0.4pt, fill=gray!50] coordinates {
        (vadd,20.5) (gelu,46.8) (softmax,60.9) (fft,52.6)
        (mmul,111.7) (attention,300.5) (conj\_grad,890.6)};
      % Эмулятор со сбором профайлинга
      \addplot[ybar, draw=black, line width=0.4pt, fill=gray!90] coordinates {
        (vadd,34.8) (gelu,121.8) (softmax,268.4) (fft,218.0)
        (mmul,970.5) (attention,2074.1) (conj\_grad,1879.7)};
      % V100-PCIE (значения переведены из мкс в мс)
      \addplot[ybar, draw=black, line width=0.4pt, fill=white] coordinates {
        (vadd,0.0072) (gelu,0.0080) (softmax,0.0078) (fft,0.0089)
        (mmul,0.0178) (attention,0.1319) (conj\_grad,0.4763)};
      \legend{Эмулятор (без проф.), Эмулятор (со сбором проф.), V100-PCIE}
    \end{axis}
  \end{tikzpicture}
  \caption{Время исполнения ядер: эмулятор без профайлинга, со сбором
           профайлинга и реальный GPU V100-PCIE. Логарифмическая шкала;
           значения V100 переведены из микросекунд в миллисекунды.}
  \label{fig:perf_comparison}
\end{figure}

% Требует в преамбуле: \usepackage{pgfplots}\pgfplotsset{compat=1.18}
\begin{figure}[h]
  \centering
  \begin{tikzpicture}
    \begin{axis}[
        width=15cm, height=8.5cm,
        ybar=1.5pt,
        bar width=8pt,
        ymode=log,
        log basis y=10,
        log origin y=infty,
        ymin=0.001, ymax=10000,
        ylabel={Время исполнения, мс (лог. шкала)},
        symbolic x coords={vadd,gelu,softmax,fft,mmul,attention,conj\_grad},
        xtick=data,
        x tick label style={font=\small},
        ymajorgrids=true,
        grid style={gray!25},
        enlarge x limits=0.08,
        legend style={at={(0.5,-0.16)}, anchor=north, legend columns=3,
                      font=\small, draw=none},
        legend cell align=left,
      ]
      % Эмулятор без профайлинга
      \addplot[ybar, draw=black, line width=0.4pt, fill=gray!50] coordinates {
        (vadd,20.5) (gelu,46.8) (softmax,60.9) (fft,52.6)
        (mmul,111.7) (attention,300.5) (conj\_grad,890.6)};
      % Эмулятор со сбором профайлинга
      \addplot[ybar, draw=black, line width=0.4pt, fill=gray!90] coordinates {
        (vadd,34.8) (gelu,121.8) (softmax,268.4) (fft,218.0)
        (mmul,970.5) (attention,2074.1) (conj\_grad,1879.7)};
      % V100-PCIE (значения переведены из мкс в мс)
      \addplot[ybar, draw=black, line width=0.4pt, fill=white] coordinates {
        (vadd,0.0072) (gelu,0.0080) (softmax,0.0078) (fft,0.0089)
        (mmul,0.0178) (attention,0.1319) (conj\_grad,0.4763)};
      \legend{Эмулятор (без проф.), Эмулятор (со сбором проф.), V100-PCIE}
    \end{axis}
  \end{tikzpicture}
  \caption{Время исполнения ядер: эмулятор без профайлинга, со сбором
           профайлинга и реальный GPU V100-PCIE. Логарифмическая шкала;
           значения V100 переведены из микросекунд в миллисекунды.}
  \label{fig:perf_comparison}
\end{figure}
